\documentclass[10pt,a4paper]{amsart}
\usepackage[margin=19mm]{geometry}
\usepackage{amsmath,amssymb,mathtools}
\usepackage[scr=rsfs,cal=euler]{mathalfa}
\usepackage{xfrac}
\usepackage[unicode,hidelinks,bookmarks=false]{hyperref}
\newcommand{\ie}{i.e.\ }
\newcommand{\cf}{cf.\ }
\newcommand{\resp}{respectively\ }
\newcommand{\Subsection}{\subsection}
\providecommand{\nobreakdash}{\nobreak\hbox{-}}
\setlength{\emergencystretch}{2em}
\multlinegap0pt
\usepackage{enumerate}
\usepackage{amssymb}
\usepackage{bm}
\usepackage{mathtools}
\newtheorem{lemma}{Lemma}
\def\G{\mathbb{G}}
\def\R{\mathbb{R}}
\newcommand{\m}{\mbox}
\title{Critical concave-convex problems in Carnot groups}
\author{Mattia Galeotti, Eugenio Vecchi}
\date{Source: arXiv:2512.04640v3 (CC BY 4.0)}
\begin{document}
\maketitle

\noindent\emph{Source and licence.} Critical concave-convex problems in Carnot groups. Version arXiv:2512.04640v3 (CC BY 4.0), distributed by arXiv under the Creative Commons Attribution 4.0 International licence, http://creativecommons.org/licenses/by/4.0/, as recorded in the arXiv licence metadata for this exact version and shown on the abstract page https://arxiv.org/abs/2512.04640v3. Full licence notice: \texttt{licenses/arxiv-2512.04640-CC-BY-4.0.txt}.\par\smallskip

\noindent\textit{Editorial scope.} Counted unit: the article's lemma lem:LambdaInf (source0.tex lines 1208--1210). The article's complete original proof of it is delivered with it (source0.tex lines 1211--1228, one \texttt{proof} environment of 18 lines). No mathematics is written, repaired or re-derived: the statement and the proof are the article's own lines, in the article's own notation, and no retained line is substituted or re-flowed.  Retained uncounted as the source-local context the proof needs: lines 264--270; lines 1162--1164.  Nothing was omitted from any retained span, so the package carries no omission record.  No retained line contains a citation, so the wrapper carries no bibliography and no bibliography entry.  No retained line contains a figure environment, an image-inclusion command or drawing code, so no image is delivered and the package declares no asset dependency.  Editorial formatting only, in addition: an AMS \texttt{amsart} preamble that restates the article's own declarations and macros used by the retained lines, namely the environments \texttt{lemma} on separate counters, together with the standard packages those lines need.  The retained environments print in this document as Lemma 1 in order of appearance, whereas the article prints the same items with the numbering of its own full document.  The complete native source of the article as distributed in the arXiv e-print for this exact version is bundled unchanged as \texttt{sources/190343/source0.tex}.
\begin{equation}\tag{{$\mathrm{P}_{\lambda}$}}\label{EqProblem}
	\left\{\begin{array}{rll}
		-\Delta_{\mathbb{G}}u& = \lambda\, u^{q} + u^{2^{\star}_{Q}-1} & \textrm{ in } \Omega,\\
		u&>0 & \textrm{ in } \Omega,\\
		u&=0 & \textrm{ on } \partial \Omega.
	\end{array}\right.
\end{equation}

\begin{equation}\label{eq:DefinitionLambda}
	\Lambda := \sup \{ \lambda >0: \eqref{EqProblem} \textrm{ admits a weak solution}\}.
\end{equation}

\begin{lemma}\label{lem:LambdaInf}
	Let $\Lambda$ as defined in \eqref{eq:DefinitionLambda}. Then, $\Lambda < +\infty$.
\end{lemma}

	\begin{proof}
	We consider the first eigenfunction $e_1$ of the operator $-\Delta_\G$
	with respect to the first Dirichlet eigenvalue $\mu_1$. In particular, the following characterization holds:
	\[
	\mu_1=\min\left\{\| |\nabla_\G u|\|^2_{L^2(\Omega)}:\ u\in S^1_0(\Omega)\m{ and }\|u\|^2_{L^2(\Omega)}=1\right\}.
	\]
	Knowing that $\|e_1\|_{L^2}=1$, $e_1>0$ a.e.~in $\Omega$ and $\| |\nabla_\G e_1|\|^2_{L^2}=\mu_1$,
	we have that any solution $u$ to \eqref{EqProblem} for some $\lambda$
	must verify
	\[
	\int_\Omega \langle \nabla_\G u,\nabla_\G e_1\rangle =\mu_1 \int_\Omega u e_1=\int_\Omega \lambda u^qe_1+u^pe_1.
	\]
	As $q\in (0,1)$ and $p>1$, for $\Lambda^*$ sufficiently big, we have
	\[
	\Lambda^*x^q+x^p>\mu_1x\quad\forall x\in \R^+.
	\]
	Therefore we must have $\lambda <\Lambda^*$ and this  proves $\Lambda\leq \Lambda^*<+\infty$.
\end{proof}

\end{document}
